\section{Discussion and Conclusion}
\label{sec:conclusion}
% ====================================================================

\subsection{Summary}

We presented a teacher--student parkour pipeline that uses monocular
depth estimates instead of a depth camera.  The GT-depth baseline
shows that the student architecture is not the limiting factor; the
main bottleneck is depth quality.  Larger DA2 models and structured
terrain texture close most of the gap to the GT-depth baseline.

The student tolerates moderate lighting and color shifts, but fails on
featureless surfaces and some unseen textures.
DSQ is a useful indicator of whether the visual pipeline preserves
terrain structure, but it does not fully predict policy success.
Overall, monocular depth appears to be a plausible bridge from RGB to
locomotion, but actual sim-to-real transfer remains unvalidated.

\subsection{Limitations}

All experiments are in simulation, so hardware transfer is still
untested.  The task is also narrow: the obstacles are simpler than
those in recent real-world parkour systems, and the robustness study
covers only two difficulty levels.

DA2 degrades on weak-texture terrain, and DSQ does not explain every failure on unseen textures.
Because DA2 processes frames independently, its estimates may flicker.
We did not measure temporal consistency or test video-based depth
models.
Deployment would also introduce failure modes absent here, including
transparent or reflective surfaces, thin structures, motion blur, and
close-range objects near the edge of the training distribution.
Finally, the perception stack is still too heavy for onboard
deployment.  The current system is best viewed as a proof of concept.

\subsection{Future Work}

The next step is hardware evaluation on the Unitree Go2 under real
camera noise, motion blur, exposure shifts, and outdoor lighting.
Broader evaluation on harder courses would test generality.  Smaller
or temporally consistent depth models could reduce compute and
stabilize predictions, while joint adaptation of the depth estimator
and student policy may improve robustness on harder textures.
